\section{Controlli supplementari del problema inverso}\label{app:inverse}
La Tabella~\ref{tab:app-excitation} esplicita la transizione tra flusso nullo e condizioni attive. Il fattore $\kappa_q$ scala simultaneamente la forza di volume e la sorgente fluida. Il numero di P\'eclet globale \`e la diagnostica $\mathrm{Pe}_{\rm glob}=\|\mathbf q_h\|_{L^\infty}L_{\rm char}/(\phi_0D)$, con $L_{\rm char}=1$; non viene utilizzato come parametro esclusivo della risposta.
\begin{table}[H]
\centering\caption{Eccitazione, rango e condizionamento del Jacobiano normalizzato sul quadrato.}\label{tab:app-excitation}
\small\setlength{\tabcolsep}{3pt}
\begin{tabular}{@{}rrrr@{}}\toprule
$\kappa_q$&$\mathrm{Pe}_{\rm glob}$&rango&$\kappa_2(J)$\\\midrule
0&0&2&--\\
$10^{-3}$&$0.00359$&4&$26571$\\
$10^{-2}$&$0.0359$&4&$2649.6$\\
$0.10$&$0.3585$&4&$258.1$\\
$0.50$&$1.7927$&4&$48.0$\\
1&$3.5854$&4&$24.78$\\
2&$7.1708$&4&$21.38$\\\bottomrule
\end{tabular}
\end{table}

\begin{table}[H]
\centering\caption{Mediane delle stime CPSO per cinque realizzazioni a ciascun livello di rumore sintetico. Le due campagne usano differenti rapporti dati/prior.}\label{tab:app-noise}
\small\setlength{\tabcolsep}{4pt}
\begin{tabular}{@{}lrrr@{}}\toprule
Parametro&Nominale&$\eta=0.05$&$\eta=0.10$\\\midrule
$\alpha$&0.700&0.7027&0.6847\\
$\beta$&0.450&0.4020&0.4234\\
$D$&0.080&0.08064&0.07587\\
$E$&2.500&2.6127&2.2282\\\bottomrule
\end{tabular}
\end{table}

La diagnostica di recupero in Figura~\ref{fig:app-recovery} mostra le realizzazioni individuali e il valore nominale. Per il caso non perturbato si usa $W_A$ e $\lambda_A=10^{-4}$; nelle campagne rumorose $W_{B,\rm cov}$ con $\lambda_B=7.2$. Alla perturbazione $\eta=0.10$ lo stesso rapporto tra pesatura dati e regolarizzazione del livello $0.05$ corrisponderebbe a $\lambda_B=1.8$: la separazione dei protocolli \`e necessaria per leggere correttamente la dispersione delle stime. Il confronto secondario con Differential Evolution e Powell utilizza un budget massimo di $600$ valutazioni del modello diretto, senza modificare il problema inverso definito in \eqref{eq:objective}.
\begin{figure}[H]
\centering\includegraphics[width=\columnwidth]{09_recovery.pdf}
\caption{Stime individuali di $(\alpha,\beta,D,E)$ nei twin experiments, per dati puliti e per due livelli di perturbazione sintetica; linee orizzontali ai valori nominali.}\label{fig:app-recovery}
\end{figure}


